% ===== 原文 PDF pp.20--23（印刷 pp.19--22）=====
% 译者：tr-02；tag 区间 2.28--2.37
% [边界] p.20 首行「and equivalently $D(\Phi,q,\tilde q)=F(\Phi,q,\tilde q)=0$.^{11}」属 p.19
% 上开始的跨页句，由 tr-01 写完（含脚注 11）；本片段自该句之后的「To be concrete, ...」续写，不重译上文。

% [原文 p.19]
具体起见，下面我们考虑规范群为 $G=\SU(N)$、含 $N_F$ 个基础表示物质的 SQCD，物质记为
$H^A_{\cN=2}$，$A=1,\dots,N_f$；不过我们写下的方程对一般的 $\cN=2$ 拉格朗日理论都成立，
只需对指标作相应改动。$D$ 项取如下形式，
\begin{equation}
  D=\frac{1}{g^2}[\Phi,\Phi^\dagger]+\big(q_A(q^A)^\dagger-(\tilde q_A)^\dagger \tilde q^{A}\big)\big|_{\text{traceless}}=0\,,
  \tag{2.28}
\end{equation}
而 $F$ 项为
\begin{equation}
  \begin{aligned}
    F_\Phi&=q_A\tilde q^{A}\big|_{\text{traceless}}=0\,,\\
    F_q&=\Phi q_A=0\,,\\
    F_{\tilde q}&=\tilde q^{A}\Phi=0\,.
  \end{aligned}
  \tag{2.29}
\end{equation}
这里 $\Phi$ 处于规范群 $\SU(N)$ 的伴随表示，$q_A$（$\tilde q^{A}$）记 $H^A_{\cN=2}$ 的标量分量，
它们分别按规范群 $\SU(N)$ 与味群 $\SU(N_F)$ 的基础（反基础）表示变换。式 (2.28) 与 (2.29)
中的关系可化简为下面几组方程：
\begin{equation}
  [\Phi,\Phi^\dagger]=0\,,
  \tag{2.30}
\end{equation}
\begin{equation}
  \left\{\begin{matrix}
    \Big(q_A(q^A)^\dagger-(\tilde q_A)^\dagger \tilde q^{A}\Big)\big|_{\text{traceless}}&=0\,,\\[2pt]
    q_A\tilde q^{A}\big|_{\text{traceless}}&=0\,,
  \end{matrix}\right.
  \tag{2.31}
\end{equation}
\begin{equation}
  \left\{\begin{matrix}
    \Phi q_A=\tilde q^{A}\Phi^\dagger&=0\,,\\[2pt]
    \Phi^\dagger q_A=\tilde q^{A}\Phi&=0\,.
  \end{matrix}\right.
  \tag{2.32}
\end{equation}

\begin{exercisebox}{习题 2.2}
  把式 (2.28) 与 (2.29) 化成式 (2.30) 至 (2.32) 的形式。
\end{exercisebox}

应当注意，式 (2.30) 至 (2.32) 具有如下特征：它们被组织成三个 $\fsu(2)_R$ 多重态，
这保证了其解保持完整的 $\cN=2$ 超对称。此外，式 (2.30) 至 (2.32) 的解有三类分支
（其示意见图 2），分别对应库仑分支、希格斯分支与混合分支，取决于哪些 R 对称性子群被破缺。

% [Tutor Note: (2.28)/(2.31) 的第一项是 $q_A(q^A)^\dagger-(\tilde q_A)^\dagger\tilde q^A$：$q_A$ 带下指标（属基础表示）、$q^A$ 带上指标（属反基础表示），两者互为厄米共轭，故这一项自动是规范群伴随表示中的一个矩阵（与 $[\Phi,\Phi^\dagger]$ 同型）；第二项 $q_A\tilde q^A$ 是两个不同手征超场之间的双线性，也处于伴随表示。原文用 "traceless" 下标强调只取无迹部分，因为整体 $U(1)$ 部分由 $D$ 项的迹部分另行处理。另注意原文 (2.28) 的排印「没有」外层方括号（1600 dpi 逐字符核对：全式只有一对 $\big[\Phi,\Phi^\dagger\big]$ 的方括号与一对 $\big(\,\big)$ 圆括号），字面上无迹竖线只挂在物质项上；物理上 $D$ 项应对整个和取无迹（$1/g^2$ 也应乘在整个和上），译文照原文排印，读者按整体无迹理解，对照 (2.31) 第一行可确认 "traceless" 的作用对象。]
% [Tutor Note: 原文在 p.20 同一页内把味指标的数目写成两种大小写：「$N_F$ 个基础表示物质」与「$A=1,\dots,N_f$」，后面又称味群为 $\SU(N_F)$（600 dpi 裁图复核 p.20 第 1 段与第 3 段，两处分别为 $N_F$、$N_f$）。按上下文应为同一个量（味指标的数目），译文照原文保留两种写法，提醒读者不要误以为是两个不同的量。]

% [原文 p.20]
\begin{definitionbox}{定义 2.3：库仑分支 $\cM_{\mathrm{CB}}$}
  库仑分支（Coulomb Branch, CB）$\cM_{\mathrm{CB}}$ 指的是超多重态标量 $q$ 与 $\tilde q$ 的真空期望值
  （VEV）均为零的情形。此时唯一需要施加的非平凡方程是 (2.30)，它要求矢量多重态标量
  $\Phi$ 在规范变换下取值于规范代数 $\fg$ 的 Cartan 子代数之中。相应地，$\cM_{\mathrm{CB}}$
  的复维数为
  \begin{equation}
    \dim_{\C}(\cM_{\mathrm{CB}})=\rank(G)\,.
    \tag{2.33}
  \end{equation}
  用规范不变算子来说，$\cM_{\mathrm{CB}}$ 由定义 2.1 中引入的 CB 算子的 VEV 来参数化。\footnote{%
  注意这些算子在有或没有共形对称性时都有意义。事实上，在非共形理论中，CB 手征主算子仍然是
  BPS 的，满足同样的缩短条件。} 由于这些算子携带非平凡的 $\mathrm{U}(1)_r$ 荷、却是
  $\SU(2)_R$ 单态，故沿 $\cM_{\mathrm{CB}}$ 方向 $\mathrm{U}(1)_r$ 被破缺，而 $\SU(2)_R$ 得以保持。
\end{definitionbox}

例如，当规范群为 $G=\SU(N)$ 时，独立的手征 CB 算子是 $\Tr(\Phi^k)$，$k=2,\dots,N$，
反手征的那些则由 $\Tr(\bar\Phi^k)$ 给出。此外，$\Phi$ 的非平凡 VEV 把 UV 规范群 $G$ 破缺到
极大环面子群 $\mathrm{U}(1)^r$，其中 $r$ 是 $G$ 的秩；这表明该分支上的低能有效描述是一个
$\mathrm{U}(1)^r$ 的 4d $\cN=2$ 超对称规范理论。

注意以上讨论迄今仍是经典的。一般而言，可以预期量子效应会修正这些解以及相应的低能 EFT。
的确，在 CB 上产生平坦度量的经典 Kähler 势会收到非平凡的量子修正，从而带来丰富的物理。
通常，QFT 中的这类量子效应（尤其是来自非微扰起源的那些）是出了名地难以研究。不过，
$\cN=2$ 超对称对量子 CB 施加了大量限制：后者必须是一类特殊的 Kähler 流形。这正是第 3 节的主题。

% [Tutor Note: 定义 2.3 中「只有 (2.30) 是非平凡的」的理由：$q=\tilde q=0$ 时 (2.31)(2.32) 自动满足，只剩 $[\Phi,\Phi^\dagger]=0$；而该式把 $\Phi$ 限制在 Cartan 子代数（差一个规范变换），故 CB 的复维数就是 $\rank G$，与 (2.33) 一致。接下来那句「把 $G$ 破缺到极大环面 $\mathrm{U}(1)^r$」与 (2.33) 是同一件事的两种说法：$r=\rank G$ 既是 $\cM_{\mathrm{CB}}$ 的复维数，也是低能剩余规范群的个数。]

\begin{definitionbox}{定义 2.4：希格斯分支 $\cM_{\mathrm{HB}}$}
  模空间的希格斯分支（Higgs Branch, HB）是指矢量多重态标量 $\Phi$ 的 VEV 为零、而超多重态
  标量 $q$ 与 $\tilde q$ 不再为零的情形。此时必须求解的非平凡方程是式 (2.31) 中的那些：
  \begin{equation}
    \cM_{\mathrm{HB}}=\left\{\begin{matrix}
      (q_a q^{\dagger b}-\tilde q_a^{\dagger}\tilde q^{b})(T^i)_b^{a}=0\\[2pt]
      (q_a\tilde q^{b})(T^i)_b^{a}=0
    \end{matrix}\right\}\Big/\{G-\text{gauge transformation}\}\,,
    \tag{2.34}
  \end{equation}
  % [原文 p.21]
  这使 HB 成为一个超 Kähler 流形。\footnote{对 $n_H$ 个超多重态，这定义了一个超 Kähler 商
  $\C^{2n_H}///G$，其矩映射由 (2.34) 第一个花括号内的方程给出。} 这里我们显式写出了规范群
  指标 $a,b$，并省略了逐对缩并的味指标；$T^i$ 是规范群 $G$ 的生成元。HB 的复维数为
  \begin{equation}
    \dim_{\C}(\cM_{\mathrm{HB}})=2(n_H-n_V)\,,
    \tag{2.35}
  \end{equation}
  其中 $n_H$ 与 $n_V$ 分别是参与希格斯化的超多重态与矢量多重态的数目。其 VEV 参数化 HB 的
  规范不变算子就是定义 2.2 中的那些算子，它们一般满足非平凡的手征环关系。\footnote{%
  再次提醒，这些算子在非共形背景下仍可看作 BPS 算子。} HB 算子在 $\mathrm{U}(1)_r$ 下中性、
  在 $\SU(2)_R$ 下带荷。相应地，HB 保持 $\mathrm{U}(1)_r$ 对称性，而破缺 $\SU(2)_R$
  对称性（以及味对称性）。
\end{definitionbox}

% [Tutor Note: (2.35) 里的 $n_H,n_V$ 是「参与希格斯化」的那部分场数（原文 `that participate in the Higgsing`），不是整个理论的场数：若某个规范因子完全没有被希格斯化，就不能把它计入 $n_V$。这与脚注 13 的超 Kähler 商表述相呼应——商掉的是真正被 Higgsing 掉的那部分规范群。]
% [Tutor Note: 定义 2.3 与 2.4 是一对对偶的破缺模式：CB 上 $\SU(2)_R$ 保住、$\mathrm{U}(1)_r$ 破缺（因为 CB 算子 $\Tr\Phi^k$ 只带 $\mathrm{U}(1)_r$ 荷）；HB 上恰好相反，$\mathrm{U}(1)_r$ 保住、$\SU(2)_R$（以及味对称性）破缺（因为 HB 算子由 $q,\tilde q$ 的双线性构成，在 $\SU(2)_R$ 下带荷）。原文两句「Correspondingly」正是这个对照。]

希格斯分支上的低能有效理论由靶空间为 $\cM_{\mathrm{HB}}$ 的 $\cN=2$ 超对称 sigma 模型描述。
与 CB 相反，由于 SUSY 非重整化定理 [117]，HB 上的超 Kähler 度量不会收到量子修正。

\begin{definitionbox}{定义 2.5：混合分支 $\cM_{\mathrm{MB}}$}
  最后一种可能是所谓的混合分支：此时矢量多重态标量与超多重态标量同时被打开，并受式 (2.30)
  至 (2.32) 的约束。用规范不变算子来说，定义 2.1 与 2.2 中的两类算子都有非零 VEV。相应地，
  在 $\cM_{\mathrm{MB}}$ 的一般点上，$\mathrm{U}(1)_r$ 与 $\SU(2)_R$ 两个 R 对称性都被破缺。
\end{definitionbox}

混合分支服从与希格斯分支类似的非重整化定理 [117]：它在局部是一个特殊 Kähler 底流形与一个
超 Kähler 纤维的度量直积。因此，相应的 EFT 是上面所述 CB 与 HB 两种 EFT 的杂交。

$\cM_{\mathrm{MB}}$ 的一个简单例子是 $\cN=4$ SYM，它可以看作 $\cN=2$ SYM 与一个处于规范群
$G$ 的伴随表示的超多重态耦合。该理论完整的模空间在度量意义上是 $\C^{3r}$，由位于 $G$ 的
Cartan 子代数中的 $\Phi$、$q$ 与 $\tilde q$ 参数化；这是一个混合分支（在此情形下实际上是
一个增强的库仑分支 [118]）。关于 $\cN=2$ 理论中混合分支的更多讨论，读者可参阅 [63,118]。

% [Tutor Note: 定义 2.5 强调「在 $\cM_{\mathrm{MB}}$ 的一般点上」（at a generic point）两个 R 对称性都破缺；这与紧接着的 $\cN=4$ SYM 例子并不矛盾——该例中整块模空间都算混合分支，原文只是补充说明它同时也可称作「增强的库仑分支」[118]。判断某一点属于哪类分支，依据是该点的无质量谱与哪些 R 对称性被破缺，而不是固定的理论。]

% [原文 p.22]
\SUB{2.4}{全纯性与 $\cN=2$ 重整化群流}

到目前为止，我们的讨论大多停留在经典层面。然而，正如上面所暗示的，库仑分支上的有效理论
通常会收到量子修正。因此，要想完整理解低能动力学，就必须掌握这类量子效应。为此，
重整化群（RG）流的一些基本方面会很有用。正如我们将要说明的，RG 跑动与 4d $\cN=2$
超对称相结合，会对量子修正给出惊人的限制。

魔力来自全纯性。举例来说，$\cN=1$ 非重整化定理 [119--121] 基本上是说：在微扰论中（在某个
特定的重整化方案下）超势 $W$ 不被重整化；而这在很大程度上要归功于超势的全纯性 [1]。
另一方面，$\cN=1$ 的 Kähler 势则可以以波函数重整化的形式收到非平凡的量子修正。但在
$\cN=2$ 超对称的情形下有更强的限制，因为 $\cN=2$ 超对称（更确切地说，是 $\SU(2)_R$
R 对称性）把超势与 Kähler 势捆绑在一起。于是，Kähler 势没有独立的量子修正。

对 4d $\cN=2$ 超势中的全纯变量，仍然存在可能的全纯重整化。让我们集中于 (2.23) 中的
$\cN=2$ SYM 部分，其复化规范耦合定义于 (2.24)。在一般规范理论中，规范耦合 $g$ 在微扰论中
每一圈阶都会收到量子修正，此外还有来自瞬子的非微扰贡献。特别地，微扰修正遵循重整化群方程
\begin{equation}
  E\frac{dg}{dE}=-\frac{g^3}{16\pi^2}b+\mathcal O(g^5)\,.
  \tag{2.36}
\end{equation}
这里 $E$ 是测量 $g$ 所用的能标，RHS 的第一项代表 1 圈 beta 函数，第二项来自更高圈的修正。
$b$ 称为 1 圈 beta 函数系数，可由相应规范理论中的场内容提取 [122,123]，即
\begin{equation}
  b=\frac{11}{3}T(\mathrm{adj})-\frac{2}{3}T(\rho_f)-\frac{1}{3}T(\rho_s)\,,
  \tag{2.37}
\end{equation}
其中 $\rho_f$ 与 $\rho_s$ 分别记规范理论中费米子与标量所处于的 $G$ 表示。

% [Tutor Note: (2.37) 是「一般」（不必超对称）规范理论的一圈系数的公式：$11/3\,T(\mathrm{adj})$ 来自规范场自身，$-2/3\,T(\rho_f)$ 按每个外线费米子计，$-1/3\,T(\rho_s)$ 按每个复标量计。把它用到 $\cN=2$ 的场内容（矢量多重态 = $\cN=1$ 矢量 + 伴随 $\cN=1$ 手征；超多重态 = 两个 $\cN=1$ 手征）就会合并成 (2.41) 的 $b=2(T(\mathrm{adj})-T(\rho_{\text{matter}}))$。可用 $\SU(N)$、$T(\mathrm{adj})=N$、$T(\text{基础})=\tfrac12$ 自检：$b=2N-N_f$，正是教科书里 $\cN=2$ SQCD 的结果。注意两式的 $T$ 归一化必须一致——见脚注 15：$T(\mathrm{adj})$ 取为 $G$ 的对偶 Coxeter 数。]
